\chapter*{Abstract}
\addcontentsline{toc}{chapter}{Abstract}

Modern spacecraft increasingly rely on commercial off-the-shelf (COTS) hardware, reducing development cost but expanding the attack surface. The bootloader is the most exposed component: it executes before any other software and, if compromised, enables persistent unauthorised code execution for the entire mission lifetime. Despite this risk, Secure Boot has not seen broad adoption in space, primarily due to stringent reliability requirements, limited resources, and the high cost of qualifying new technology for space.

This thesis investigates how Secure Boot can be incorporated into space mission bootloaders without compromising these constraints. A lifecycle-based threat model, grounded in the ENISA Space Threat Landscape taxonomy, identifies two primary attacker scenarios: a physical adversary with pre-launch hardware access, and a remote adversary exploiting firmware update channels after deployment. These scenarios motivate a formal Secure Boot specification, designated SAVOIR-GS-002S, written in the language and requirement structure of the ESA SAVOIR framework so that it can be adopted by standards bodies without a translation step. The specification introduces seven additive security requirements covering a hardware Root of Trust, digital signature verification, non-bypassability, protection against both classical and quantum computational attacks, rollback protection via a monotonic counter, and structured failure reporting.

A hardware trade-off analysis shows that only two primitives are strictly necessary to realise the specification: sector-level hardware-enforced write protection for the Boot Software and public key storage, and protected non-volatile storage for the monotonic rollback counter. Both are present on commodity ARM Cortex-M microcontrollers without requiring any specialised hardware. A proof-of-concept implementation in C on the STM32F439ZI confirms that all seven requirements are achievable on representative COTS hardware. The implementation includes A/B image partitioning, monotonic rollback protection in option-byte-protected flash, PUS-compatible boot reporting, and an algorithm-agnostic cryptographic verification interface, with an automated end-to-end test suite verifying all nominal and failure paths over UART and SWD.

Performance benchmarks for ECDSA-P256, RSA-3072/4096, ML-DSA-44/65, and LMS demonstrate that the security overhead is modest. A hybrid ECDSA-P256 plus LMS deployment on a 128~kB image completes in approximately 363~ms with hardware SHA-2 acceleration, well within a one-second boot budget, and LMS requires only 3~kB of working memory and a 60-byte public key. The results show that incorporating Secure Boot into space-grade bootloaders is not a fundamental hardware or performance barrier, but primarily a gap in standards and tooling. The SAVOIR-GS-002S specification and accompanying implementation provide a concrete, standards-compatible starting point for closing that gap.